%=============================================================================
% PATENT 15 — NEW PATENT PROPOSALS
% Specific apparatus configurations and measurement techniques
% identified from the LSDFG engineering design
%=============================================================================
\documentclass[11pt,a4paper]{article}

\usepackage[margin=1in]{geometry}
\usepackage{amsmath,amssymb}
\usepackage{siunitx}
\usepackage{booktabs}
\usepackage{longtable}
\usepackage{hyperref}
\usepackage{xcolor}
\usepackage{enumitem}
\usepackage{tcolorbox}

\title{\textbf{New Patent Proposals Arising from Patent 15}\\[6pt]
\large Laboratory Scalar Density Field Generator\\
\normalsize Technology Opportunities Identified During
Engineering Design}

\author{DFD Research Programme}
\date{March 2026}

\begin{document}
\maketitle
\tableofcontents
\newpage

%=============================================================================
\section{Overview}
%=============================================================================

During the detailed engineering design of the Laboratory Scalar
Density Field Generator (Patent 15, US 63/868,073), four specific
apparatus configurations and measurement techniques were identified
that merit independent patent protection.  Each represents a
distinct inventive step beyond the master methods patent and
addresses a specific application domain.

%=============================================================================
\section{Proposal A: Dual-Mode Cavity $\psi$-Pump with
Phase-Controlled Geometry Breaking}
%=============================================================================

\subsection{Title}

\textit{``Apparatus and Method for Scalar Field Generation Using
Simultaneously Excited TE and TM Modes in a Superconducting
Resonant Cavity with Controlled Phase Offset''}

\subsection{Problem Addressed}

The geometry transparency property of symmetric resonant cavities
operating in pure eigenmodes causes the EM-to-$\psi$ driving
overlap integral $G$ to vanish identically.  This means that
conventional cavity operation cannot generate $\psi$ perturbations
through the driven-resonance channel, regardless of the value of
$|\lambda - 1|$.  Breaking this transparency requires a specific
apparatus configuration that has not been previously disclosed.

\subsection{Inventive Step}

The invention is a resonant cavity apparatus configured to
simultaneously sustain TE$_{0mn}$ and TM$_{0m'n'}$ modes with a
controlled relative phase $\phi$, such that the time-averaged
electromagnetic stress tensor trace $\Xi = -(1/2)(B^2 - E^2/c^2)$
acquires a nonzero spatial integral weighted against the target
$\psi$-mode profile.

Key elements:
\begin{enumerate}[nosep]
\item \textbf{Dual-port excitation}: Two independent RF feeds
at frequencies $\omega_{\rm TM}$ and $\omega_{\rm TE}$,
phase-locked with programmable offset $\phi \in [0, 2\pi)$.
\item \textbf{Phase-controlled overlap restoration}: The driven
overlap is $G(\phi) = u(z_0)\,e^{-2\psi_0}\,\eta_\times\,U_0
\cos\phi$, maximized at $\phi = 0$ or $\pi$.
\item \textbf{Cross-mode overlap coefficient}: The parameter
$\eta_\times$ depends on the specific mode pair and cavity
geometry; the invention includes optimized cavity geometries
that maximize $\eta_\times$.
\item \textbf{Frequency matching}: Cavity geometry designed so
that $\omega_{\rm TM} + \omega_{\rm TE} = \Omega_\psi$
(target $\psi$-mode frequency) or $2\omega_{\rm TM} =
\Omega_\psi$ with TE providing the spatial symmetry breaking.
\item \textbf{Phase modulation}: Deliberate $2\omega$
amplitude modulation at depth $m$ to enable both driven
and parametric pumping channels simultaneously.
\end{enumerate}

\subsection{Claims Outline}

\begin{enumerate}
\item An apparatus for generating scalar density field
perturbations comprising a resonant cavity simultaneously
excited in at least one TE mode and at least one TM mode
with a controlled relative phase offset.

\item The apparatus of claim 1 wherein the resonant cavity
is a superconducting RF cavity with $Q_0 > 10^8$ and the
relative phase is controlled to within $0.01\;\si{rad}$.

\item The apparatus of claim 1 further comprising a
phase-locked loop maintaining the frequency sum or
difference of the two modes at a target $\psi$-mode
frequency $\Omega_\psi$.

\item A method for generating scalar density field perturbations
comprising: (a) exciting a TE mode and a TM mode simultaneously
in a resonant cavity; (b) controlling the relative phase to
maximize the driving overlap integral $G$; (c) applying
amplitude modulation at frequency $2\omega$ to enable
parametric amplification.

\item The method of claim 4 further comprising scanning the
target frequency $\Omega_\psi$ to identify $\psi$-mode
resonances.

\item The apparatus of claim 1 wherein the cavity geometry
is asymmetric with at least one iris or constriction to break
the equipartition condition that causes $G = 0$ in symmetric
geometries.

\item An array of cavities as described in claim 1, positioned
at antinodes of a target $\psi$ mode, with phases controlled
to constructively drive the mode.
\end{enumerate}

\subsection{Distinction from Patent 15}

Patent 15 discloses the general concept of a chamber with
configurable sources for $\psi$ generation.  This proposal
covers the specific dual-mode excitation technique, the
geometry-breaking mechanism, and the phase-control apparatus
that make $\psi$ generation physically possible through the
driven-resonance channel.

\subsection{Market Applications}

\begin{itemize}[nosep]
\item Fundamental physics: First laboratory test of
EM--$\psi$ coupling
\item Precision metrology: New probe of vacuum constitutive
relations
\item Particle accelerator facilities: Leverages existing SRF
infrastructure
\item Defense: Anomalous force generation if $|\lambda - 1|$
is sufficiently large
\end{itemize}

%=============================================================================
\section{Proposal B: Multi-Modal $\psi$-Field Sensor Fusion
System with Kalman-Filter State Estimation}
%=============================================================================

\subsection{Title}

\textit{``System and Method for Real-Time Estimation of a Scalar
Density Field Using Fused Measurements from Optical, Atomic, and
Inertial Sensors with Extended Kalman Filter State Estimation''}

\subsection{Problem Addressed}

The scalar density field $\psi$ couples to three distinct physical
domains (optical, atomic, inertial), each with different noise
characteristics, bandwidths, and systematic error profiles.  No
existing system combines these three measurement modalities into
a single coherent field estimate, nor does any existing estimator
incorporate the $\psi$ field equation as a process model.

\subsection{Inventive Step}

The invention is a sensor fusion system that:

\begin{enumerate}[nosep]
\item Combines measurements from at least two of the three
domains (optical interferometry, atomic clock comparison,
inertial acceleration) into a single estimate of
$\psi(\mathbf{x},t)$.

\item Uses the DFD field equation as the process model in an
extended Kalman filter, so that the estimate is constrained
by the known physics of $\psi$ propagation and coupling.

\item Provides real-time field estimates at rates $> 1\;\si{kHz}$,
enabling feedback control of $\psi$-field sources.

\item Automatically detects and flags sensor failures or
systematic errors through cross-validation between domains
(a $\psi$ perturbation seen optically but not inertially
indicates a systematic, not a real field).
\end{enumerate}

Key technical elements:
\begin{itemize}[nosep]
\item \textbf{Process model}: Discretized $\psi$ field equation
with screening, implemented as a sparse matrix on FPGA.
\item \textbf{Measurement model}: Domain-specific sensor response
functions (optical: path-integrated $\psi$; atomic: local
coupling coefficient; inertial: $\psi$ gradient).
\item \textbf{Adaptive noise estimation}: Online estimation of
sensor noise covariances using innovation monitoring.
\item \textbf{Multi-rate fusion}: Fast channels (optical,
inertial at $\si{kHz}$--$\si{MHz}$) and slow channels
(atomic at $\sim 1\;\si{Hz}$) fused in a multi-rate filter
architecture.
\end{itemize}

\subsection{Claims Outline}

\begin{enumerate}
\item A system for estimating a scalar density field comprising:
(a) at least one optical sensor measuring a quantity dependent
on the local refractive index $n = e^\psi$; (b) at least one
inertial sensor measuring acceleration proportional to
$\nabla\psi$; (c) a processor implementing an extended Kalman
filter with the $\psi$ field equation as process model; (d) an
output providing a real-time estimate of $\psi(\mathbf{x},t)$.

\item The system of claim 1 further comprising at least one
atomic frequency reference providing a measurement dependent
on the local value of $\psi$ through species-specific coupling
coefficients.

\item The system of claim 1 wherein the extended Kalman filter
operates at a rate exceeding $1\;\si{kHz}$ using FPGA-implemented
sparse matrix operations.

\item A method for detecting scalar density field perturbations
comprising: (a) measuring optical phase shifts in an
interferometer; (b) measuring acceleration with an inertial
sensor; (c) fusing measurements using a model-based estimator
incorporating the field equation; (d) declaring a detection
when the estimated field exceeds a threshold in at least two
independent measurement domains.

\item The system of claim 1 further comprising a cross-validation
module that flags events detected in only one domain as
potential systematic errors.

\item The system of claim 1 wherein the output estimate is
used as feedback to control sources of the scalar density field.
\end{enumerate}

\subsection{Distinction from Patent 15}

Patent 15 discloses a sensor suite and control electronics in
general terms.  This proposal covers the specific estimation
algorithm, the multi-rate fusion architecture, the cross-validation
logic, and the use of the $\psi$ field equation as a physics-based
process model---a fundamentally new approach to gravitational
field estimation.

\subsection{Market Applications}

\begin{itemize}[nosep]
\item Gravitational wave detection: Enhanced sensitivity through
multi-modal fusion
\item Precision geodesy: Centimeter-level geoid determination
\item Navigation: GPS-denied inertial navigation with
gravitational correction
\item Seismology: Separation of tectonic and gravitational
signals
\item Underground surveying: Density mapping through
$\psi$-field tomography
\end{itemize}

%=============================================================================
\section{Proposal C: Frequency-Scanning $\psi$-Mode Spectroscopy
Apparatus with Lock-In Detection}
%=============================================================================

\subsection{Title}

\textit{``Apparatus and Method for Spectroscopic Detection of
Scalar Density Field Resonances Using Swept-Frequency
Electromagnetic Excitation with Synchronous Detection''}

\subsection{Problem Addressed}

The natural frequency spectrum of $\psi$ modes in a laboratory
volume is unknown a priori.  Unlike electromagnetic cavity modes
(calculable from Maxwell's equations and boundary conditions),
$\psi$-mode frequencies depend on the unknown propagation speed
$c_s$ and the boundary conditions imposed by the local
gravitational environment (Earth's screening).  A systematic
method for discovering these modes is needed.

\subsection{Inventive Step}

The invention is a spectroscopic method and apparatus that:

\begin{enumerate}[nosep]
\item Scans the EM drive frequency $\omega$ such that $2\omega$
sweeps through the expected $\psi$-mode frequency range.

\item Uses synchronous (lock-in) detection at $2\omega$ on all
sensor channels to extract the $\psi$ response at each
frequency step.

\item Constructs a ``$\psi$-mode spectrum'' analogous to an NMR
or microwave absorption spectrum.

\item Identifies $\psi$-mode candidates as peaks in the spectrum
that appear consistently across multiple independent sensor
channels.

\item Provides automatic mode characterization (frequency,
linewidth, coupling strength) through Lorentzian fitting.
\end{enumerate}

Key technical elements:
\begin{itemize}[nosep]
\item \textbf{Frequency scanning}: Step size $\Delta f < \gamma_\psi$
(mode linewidth) to avoid missing narrow resonances.
\item \textbf{Dwell time}: At least $5/\gamma_\psi$ per step
for steady-state response.
\item \textbf{Lock-in detection}: Digital lock-in at $2\omega$
with bandwidth $\sim \gamma_\psi/10$ for optimal SNR.
\item \textbf{Multi-channel correlation}: Simultaneous lock-in
on optical, inertial, and (if available) atomic channels.
\item \textbf{Candidate ranking}: Signal-to-noise ratio, number
of confirming channels, Lorentzian fit quality.
\item \textbf{False-positive rejection}: On/off modulation of
the EM source; true $\psi$ responses must correlate with
source state.
\end{itemize}

\subsection{Claims Outline}

\begin{enumerate}
\item A method for detecting scalar density field resonances
comprising: (a) exciting an electromagnetic resonator at a
drive frequency $\omega$; (b) sweeping $\omega$ such that
$2\omega$ traverses a target frequency range; (c) at each
frequency step, performing synchronous detection at $2\omega$
on at least one $\psi$-sensitive sensor; (d) constructing a
spectrum of response amplitude versus $2\omega$; (e) identifying
resonance peaks in the spectrum.

\item The method of claim 1 wherein synchronous detection is
performed simultaneously on sensors in at least two distinct
physical domains.

\item The method of claim 1 further comprising verifying each
candidate resonance by: (i) toggling the EM source on and off;
(ii) confirming that the response at $2\omega$ correlates with
the source state.

\item An apparatus for $\psi$-mode spectroscopy comprising:
(a) a tunable electromagnetic resonator; (b) a frequency
synthesizer with sweep capability; (c) at least one
$\psi$-sensitive sensor; (d) a digital lock-in amplifier
operating at $2\omega$; (e) a processor constructing and
analyzing the response spectrum.

\item The method of claim 1 wherein the frequency scan
is adaptive, dwelling longer at frequencies where the
response exceeds a preliminary threshold.

\item The method of claim 1 further comprising characterizing
each identified resonance by fitting a Lorentzian profile to
extract the mode frequency $\Omega_\psi$, linewidth
$2\gamma_\psi$, and coupling amplitude.
\end{enumerate}

\subsection{Distinction from Patent 15}

Patent 15 discloses the concept of configurable sources and
frequency control.  This proposal covers the specific
spectroscopic measurement protocol, the lock-in detection
scheme at $2\omega$, the multi-channel correlation analysis,
and the adaptive scanning algorithm---a complete new measurement
technique for discovering $\psi$-field resonances.

\subsection{Market Applications}

\begin{itemize}[nosep]
\item Fundamental physics: Discovery of new scalar field modes
\item Materials science: Probing vacuum electromagnetic
properties in confined geometries
\item Quantum information: Identifying decoherence channels
from scalar field coupling
\item Metrology: Calibration of gravitational sensors against
known $\psi$ modes
\end{itemize}

%=============================================================================
\section{Proposal D: $\psi$-Field Safety Interlock System with
Multi-Sensor Redundant Trip Logic}
%=============================================================================

\subsection{Title}

\textit{``Safety System for Scalar Density Field Generation
Apparatus with Independent Multi-Domain Monitoring, Hardware-
Enforced Field Limits, and Microsecond-Response Shutdown
Capability''}

\subsection{Problem Addressed}

If the scalar density field can be artificially generated, the
acceleration $\mathbf{a} = (c^2/2)\nabla\psi$ means that
uncontrolled $\psi$ gradients pose a physical hazard.  A
$\psi$ gradient of $|\nabla\psi| = 10^{-16}\;\si{m^{-1}}$
produces acceleration comparable to Earth's gravity.
No existing safety system is designed to detect or limit
artificially generated gravitational fields.

\subsection{Inventive Step}

The invention is a safety interlock system specifically designed
for $\psi$-field generation apparatus, with:

\begin{enumerate}[nosep]
\item \textbf{Independent safety sensors}: Accelerometers, optical
phase monitors, and pressure sensors on a separate power rail
and data path from the measurement system.

\item \textbf{Multi-domain trip logic}: A trip is triggered when
anomalous signals are detected in at least one inertial sensor
AND confirmed by at least one optical sensor, reducing false
trips from single-sensor failures.

\item \textbf{Graduated response}: Warning level (log and alert),
power reduction level, fast RF kill ($< 1\;\si{\mu s}$), and
full shutdown ($< 1\;\si{s}$), with distinct thresholds for
each.

\item \textbf{Rate-of-change monitoring}: Trips on $d|\psi|/dt$
exceeding threshold, detecting the onset of parametric
instability before the field amplitude becomes hazardous.

\item \textbf{Hardware implementation}: Trip logic in hardwired
relay chains (not software), ensuring function even with total
computer failure.

\item \textbf{Personnel dosimetry}: Wearable accelerometer
badges logging exposure to anomalous acceleration fields,
analogous to radiation dosimetry.
\end{enumerate}

\subsection{Claims Outline}

\begin{enumerate}
\item A safety system for a scalar density field generation
apparatus comprising: (a) at least two independent inertial
sensors measuring acceleration; (b) at least one optical sensor
measuring refractive index variation; (c) hardwired trip logic
that removes power from field-generation sources when a
predefined threshold is exceeded; (d) a graduated response
with at least two distinct action levels.

\item The system of claim 1 wherein the trip logic requires
confirmation from at least two independent sensor domains
before initiating the highest shutdown level.

\item The system of claim 1 further comprising rate-of-change
monitoring that triggers a trip when the time derivative of
the estimated field amplitude exceeds a threshold.

\item The system of claim 1 wherein the RF source kill is
achieved by PIN diode switches with response time
$< 1\;\si{\mu s}$.

\item A wearable personnel safety device for use near scalar
density field generation apparatus, comprising: (a) a MEMS
accelerometer logging acceleration data; (b) a threshold alarm;
(c) a cumulative exposure integrator analogous to a radiation
dosimeter.

\item The system of claim 1 further comprising an automated
test function that periodically verifies trip response times
without generating a real $\psi$ field.
\end{enumerate}

\subsection{Distinction from Patent 15}

Patent 15 discloses safety enforcement as one function of the
control system.  This proposal covers the specific
implementation of a SIL-2 rated safety instrumented system with
hardware interlocks, multi-domain confirmation logic, personnel
dosimetry, and microsecond response---critical enabling
technology that must work correctly the first time.

\subsection{Market Applications}

\begin{itemize}[nosep]
\item Laboratory safety: Required for any $\psi$-field
generation facility
\item Regulatory compliance: Framework for government
safety certification of novel field generators
\item Personnel protection: Wearable dosimeters for
gravitational field exposure
\item Industrial safety: Extension to any high-energy
physics facility with potential anomalous force generation
\end{itemize}

%=============================================================================
\section{Patent Portfolio Strategy}
%=============================================================================

\subsection{Relationship to Existing Patents}

\begin{longtable}{p{0.12\textwidth}p{0.40\textwidth}p{0.38\textwidth}}
\toprule
\textbf{Patent} & \textbf{Coverage} & \textbf{Relation to New Proposals} \\
\midrule
\endhead
\#15 & Master methods: chamber, sources, sensors, control, safety
(general) & Parent patent; new proposals are specific embodiments \\
\#11 & EM coupling / SRF & Covers coupling physics; Proposal A covers
the specific dual-mode apparatus \\
\#6 & Quantum sensors / imaging & Covers sensor concepts; Proposal B
covers the specific fusion algorithm \\
\#5 & Navigation / timing & Covers navigation applications; Proposal C
covers the mode-discovery technique \\
New A & Dual-mode $\psi$-pump & Specific source apparatus \\
New B & Sensor fusion system & Specific measurement system \\
New C & $\psi$-mode spectroscopy & Specific measurement protocol \\
New D & Safety interlock system & Specific safety apparatus \\
\bottomrule
\end{longtable}

\subsection{Filing Priority}

\begin{enumerate}
\item \textbf{Proposal A} (Highest): The dual-mode cavity pump
is the core enabling technology.  Without it, geometry transparency
prevents $\psi$ generation.  File within 6 months to establish
priority on the specific technique.

\item \textbf{Proposal D} (High): Safety systems must be
established before any high-power operation.  The regulatory
framework for gravitational field safety does not exist; early
filing establishes the standard.

\item \textbf{Proposal B} (Medium): The sensor fusion system is
necessary for practical operation but has longer development
timelines.  File within 12 months.

\item \textbf{Proposal C} (Medium): The spectroscopy method
depends on having a working pump (Proposal A).  File within
12 months, ideally after initial experimental results.
\end{enumerate}

\subsection{Estimated Filing Costs}

\begin{longtable}{p{0.15\textwidth}p{0.25\textwidth}p{0.25\textwidth}p{0.25\textwidth}}
\toprule
\textbf{Proposal} & \textbf{Provisional (US)} & \textbf{Utility (US)}
& \textbf{PCT} \\
\midrule
\endhead
A & \$3,000 & \$15,000 & \$8,000 \\
B & \$3,000 & \$12,000 & \$8,000 \\
C & \$2,500 & \$10,000 & \$8,000 \\
D & \$2,500 & \$10,000 & \$8,000 \\
\midrule
\textbf{Total} & \$11,000 & \$47,000 & \$32,000 \\
\bottomrule
\end{longtable}

%=============================================================================
\section{Summary}
%=============================================================================

Four new patent proposals have been identified from the Patent 15
engineering design process:

\begin{tcolorbox}[colback=green!3, colframe=green!50!black,
title=New Patent Proposals Summary]
\begin{description}[nosep]
\item[Proposal A:] \textbf{Dual-Mode Cavity $\psi$-Pump} ---
The specific TE+TM superposition technique that breaks geometry
transparency and enables $\psi$ generation through the
driven-resonance channel.

\item[Proposal B:] \textbf{Multi-Modal Sensor Fusion} ---
Real-time Kalman filter estimation of $\psi(\mathbf{x},t)$
from fused optical, atomic, and inertial measurements using
the field equation as process model.

\item[Proposal C:] \textbf{$\psi$-Mode Spectroscopy} ---
Swept-frequency discovery of scalar field resonances using
$2\omega$ lock-in detection with multi-channel confirmation.

\item[Proposal D:] \textbf{Safety Interlock System} ---
Hardware-enforced $\psi$-field limits with microsecond
response, multi-domain confirmation, and personnel dosimetry.
\end{description}
\end{tcolorbox}

Together with Patent 15 (master methods), these four proposals
create a comprehensive IP portfolio covering the generation,
detection, characterization, and safe operation of laboratory
scalar density field apparatus.

\end{document}
